% In this file you should put all LaTeX macros and settings to be used both by
% the pdf version and the web version.
% This should be most of your macros.

% The theorem-like environments defined below are those that appear by default
% in the dependency graph. See the README of leanblueprint if you need help to
% customize this. 
% The configuration below use the theorem counter for all those environments
% (this is what the [theorem] arguments mean) and never resets it.
% If you want for instance to number them within chapters then you can add
% [chapter] at the end of the next line.
\newtheorem{theorem}{Theorem}
\newtheorem{proposition}[theorem]{Proposition}
\newtheorem{lemma}[theorem]{Lemma}
\newtheorem{corollary}[theorem]{Corollary}

\theoremstyle{definition}
\newtheorem{definition}[theorem]{Definition}

% Shorthand for the modulus vocabulary (families of objects, Fulkerson duality).
\newcommand{\Adm}{\operatorname{Adm}}
\newcommand{\Mod}{\operatorname{Mod}}
\newcommand{\MEO}{\operatorname{MEO}}
\newcommand{\FEU}{\operatorname{FEU}}
\newcommand{\Energy}{\mathcal{E}}

% No bibliography backend is configured, so papers are cited with a plain
% hyperlinked author/year rather than \cite. One macro per paper as they're added.
\newcommand{\feupaperref}{\href{https://doi.org/10.1016/j.disc.2020.112282}{Albin, Clemens, Hoare, Poggi-Corradini, Sit, and Tymochko (2021)}}

% Flags a place where the Lean formalization diverges from the on-paper
% definition for implementation reasons (not mathematical ones), with a link
% to the file under docs/ that explains the divergence in detail. #1 is a
% short inline summary; #2 is the docs/ filename (e.g. common.md).
\newcommand{\implnote}[2]{\par\noindent\textit{Implementation note: #1} See
  \href{https://github.com/nathan-albin/lean-modulus/blob/main/docs/#2}
  {\texttt{docs/#2}} for details.}